\hmsubsection{Computer Vision}{OAH}{}

Computer vision is a field within computer science concerned with enabling machines to interpret and understand visual information from the world. By processing digital images or video streams, a computer vision system can identify objects, determine their positions, and classify their properties automatically.

In the context of this project, computer vision is used to detect and classify coloured balls in real time. The system captures frames from a camera mounted above the workspace and analyses each frame to determine the colour and position of any balls present. This information is subsequently passed to the robot arm, which uses it to perform the physical sorting task.

The pipeline implemented in this project follows a classical structure. First, the raw image is preprocessed to reduce noise and normalise lighting conditions. Next, features are extracted from the image — in this project, colour information in the HSV colour space and geometric properties of detected contours. Finally, the extracted features are used to classify and localise objects in the scene.

A key challenge in colour-based detection is robustness to changes in ambient lighting. Variations in illumination can cause the same object to appear with significantly different colour values across frames. To address this, the system employs Contrast Limited Adaptive Histogram
Equalisation (CLAHE), a technique that enhances local image contrast while
preserving colour information. These techniques are described in detail in
Section~8.2.2 \cite{bradski2008learning}.
